\documentclass[11pt]{article}
\usepackage{mathblatt}
% gesetzt von werkzeuge/setzer.py (Sorte selbst) aus dem Lernweg E69
% und den Bankzeilen; Muster bau/proben/2026-10-09/hypotenuse-muster.
\usepackage{enumitem}
\usepackage{adjustbox,varwidth}
\newgeometry{a4paper,margin=18mm,top=15mm,bottom=20mm,footskip=9mm}
\pagestyle{fancy}\fancyhf{}\renewcommand{\headrulewidth}{0pt}
\fancyfoot[L]{\footnotesize\color{mbgrau}E69 \textperiodcentered{} Lösungen \textperiodcentered{} Kreisfläche}
\fancyfoot[R]{\footnotesize\color{mbgrau}\thepage}
\setlength{\parskip}{3pt}
\renewcommand{\kreuz}[1]{\mbox{$\square$\ #1}\quad}
\newcommand{\kreuzl}[1]{\par\hangindent1.4em\hangafter1\noindent$\square$\ #1}
\newcommand{\aufg}[3]{\par\noindent\makebox[0pt][r]{#3}\begin{minipage}[t]{\linewidth}#1\end{minipage}\par}
\newcommand{\aufgg}[4]{\par\noindent\makebox[0pt][r]{#4}\begin{minipage}[t]{0.6\linewidth}\vspace{0pt}#1\end{minipage}\hfill
  \begin{minipage}[t]{0.37\linewidth}\vspace{0pt}\centering #2\end{minipage}\par}
\newcommand{\nr}[1]{\makebox[7mm][l]{\textbf{#1}}}
\newcommand{\lsg}[3]{\par\noindent\begin{minipage}[t]{9mm}\textbf{#1}\end{minipage}%
  \begin{minipage}[t]{52mm}\raggedright\bfseries\boldmath #2\end{minipage}\hspace{3mm}%
  \begin{minipage}[t]{\dimexpr\linewidth-64mm\relax}\raggedright\small #3\end{minipage}\par\vspace{4pt}}
\newlength{\kr}
\makeatletter
\newcommand{\karomerk}[2]{\protected@write\@auxout{}{\string\karoseite{#1}{#2}{\thepage}}}
\newcommand{\karoseite}[3]{\expandafter\xdef\csname karo@#1@#2\endcsname{#3}}
\newcommand{\karohinweis}[3]{\karomerk{h}{#1}\ifcsname karo@k@#1\endcsname
  \ifnum\csname karo@k@#1\endcsname=0\csname karo@h@#1\endcsname\relax#2\else#3\fi
  \else#3\fi}
\makeatother
\begin{document}

\noindent{\Large\bfseries Lösungen}\hspace{1em}{\color{mbgrau}Kreisfläche}\par\smallskip
\noindent{\small Links das Ergebnis, rechts der Weg in kurzen Schritten. Stimmt dein Ergebnis nicht, such den ersten Schritt, der anders ist.}\par
\par\Needspace{25mm}\vspace{6pt}\noindent\colorbox{black!12}{\makebox[7mm]{\bfseries\strut A}}\hspace{2mm}{\bfseries Fläche aus Radius oder Durchmesser}\par\vspace{4pt}
\lsg{1.}{(a) 78,5\,cm² (b) 314,2\,cm² (c) 50,3\,m²}{(a) $A = \pi \cdot 25 \approx 78{,}5$\,cm²; (b) $A = \pi \cdot 100 \approx 314{,}2$\,cm²; (c) $A = \pi \cdot 16 \approx 50{,}3$\,m²}
\lsg{2.}{(a) 153,9\,cm² (b) 113,1\,cm² (c) 63,6\,m²}{(a) $r = 7$\,cm, $A = \pi \cdot 49 \approx 153{,}9$\,cm²; (b) $r = 12 : 2 = 6$\,cm, $A = \pi \cdot 36 \approx 113{,}1$\,cm²; (c) $r = 9 : 2 = 4{,}5$\,m, $A = \pi \cdot 4{,}5^2 \approx 63{,}6$\,m²}
\lsg{3.}{(a) 19,63\,cm² (b) 2,54\,m²}{(a) $A = \pi \cdot 2{,}5^2 \approx 19{,}63$\,cm²; (b) $r = 0{,}9$\,m, $A = \pi \cdot 0{,}9^2 \approx 2{,}54$\,m²}
\lsg{4.}{eine große, etwa 44\,cm² mehr}{groß: $r = 16$\,cm, $A = \pi \cdot 16^2 \approx 804$\,cm²; klein: $r = 11$\,cm, $A = \pi \cdot 11^2 \approx 380$\,cm², zwei kleine $\approx 760$\,cm²; Unterschied $\approx 44$\,cm²: Die große Pizza ist mehr.}
\par\Needspace{25mm}\vspace{6pt}\noindent\colorbox{black!12}{\makebox[7mm]{\bfseries\strut B}}\hspace{2mm}{\bfseries Halbkreis und Viertelkreis}\par\vspace{4pt}
\lsg{1.}{(a) 25,1\,cm² (b) 78,5\,cm² (c) 56,5\,m²}{(a) $A = \pi \cdot 4^2 : 2 \approx 25{,}1$\,cm²; (b) $A = \pi \cdot 10^2 : 4 \approx 78{,}5$\,cm²; (c) $r = 6$\,m, $A = \pi \cdot 6^2 : 2 \approx 56{,}5$\,m²}
\lsg{2.}{$\pi \cdot a^2 : 4$; 50,3\,cm²}{ein Viertel der Kreisfläche; für $a = 8$: $\pi \cdot 64 : 4 \approx 50{,}3$\,cm²}
\lsg{3.}{(a) 1\,963,5\,m² (b) 6\,463,5\,m²}{(a) zwei Halbkreise sind ein Kreis: $r = 25$\,m, $A = \pi \cdot 25^2 \approx 1\,963{,}5$\,m²; (b) Rechteck $90 \cdot 50 = 4\,500$\,m²; $4\,500 + 1\,963{,}5 \approx 6\,463{,}5$\,m²}
\lsg{4.}{Nein, 39,3\,m² > 35\,m²}{Nein; Halbkreis mit $r = 5$\,m: $A = \pi \cdot 5^2 : 2 \approx 39{,}3$\,m², das ist mehr als 35\,m²; es fehlen $\approx 4{,}3$\,m²}
\par\Needspace{25mm}\vspace{6pt}\noindent\colorbox{black!12}{\makebox[7mm]{\bfseries\strut C}}\hspace{2mm}{\bfseries Radius aus der Fläche}\par\vspace{4pt}
\lsg{1.}{(a) 10,0\,cm (b) 5,0\,m (c) 2,0\,cm}{(a) $r = \sqrt{314 : \pi} \approx 10{,}0$\,cm; (b) $r = \sqrt{78{,}5 : \pi} \approx 5{,}0$\,m; (c) $r = \sqrt{12{,}6 : \pi} \approx 2{,}0$\,cm}
\lsg{2.}{(a) 8,0\,cm (b) 1,6\,m}{(a) $r = \sqrt{50 : \pi} \approx 3{,}99$\,cm, $d \approx 8{,}0$\,cm; (b) $r = \sqrt{2 : \pi} \approx 0{,}80$\,m, $d \approx 1{,}6$\,m}
\lsg{3.}{$d \approx 3{,}91$\,m}{$r^2 = 12 : \pi$; $r \approx 1{,}954$\,m; $d = 2 \cdot r \approx 3{,}91$\,m}
\lsg{4.}{4\,m}{Die Ziege erreicht einen Kreis, der Strick ist der Radius. $r = \sqrt{45 : \pi} \approx 3{,}78$\,m; 3\,m sind zu kurz, 4\,m reichen.}
\par\Needspace{25mm}\vspace{6pt}\noindent\colorbox{black!12}{\makebox[7mm]{\bfseries\strut D}}\hspace{2mm}{\bfseries Fläche oder Umfang?}\par\vspace{4pt}
\lsg{1.}{(a) U (b) F (c) U (d) F (e) U}{}
\lsg{2.}{(a) 10\,cm; 31,4\,cm; 78,5\,cm² (b) 4\,m; 25,1\,m; 50,3\,m² (c) 3\,cm; 9,4\,cm; 7,1\,cm²}{(a) $d = 10$\,cm, $u \approx 31{,}4$\,cm, $A \approx 78{,}5$\,cm²; (b) $r = 4$\,m, $u \approx 25{,}1$\,m, $A \approx 50{,}3$\,m²; (c) $d = 3$\,cm, $u \approx 9{,}4$\,cm, $A \approx 7{,}1$\,cm²}
\lsg{3.}{(a) 219,9\,cm (b) 1\,963,5\,cm² (c) 282,7\,cm}{(a) Umfang: $u = \pi \cdot 70 \approx 219{,}9$\,cm; (b) Fläche: $A = \pi \cdot 25^2 \approx 1\,963{,}5$\,cm²; (c) Umfang: $u = 2 \cdot \pi \cdot 45 \approx 282{,}7$\,cm}
\lsg{4.}{38 Steine, 4 Säcke}{Steine: $u = \pi \cdot 3 \approx 9{,}42$\,m; $9{,}42 : 0{,}25 \approx 37{,}7$, aufrunden: $n = 38$ Steine; Kies: $r = 1{,}5$\,m, $A = \pi \cdot 1{,}5^2 \approx 7{,}07$\,m²; $7{,}07 : 2 \approx 3{,}53$, aufrunden: $n = 4$ Säcke}
\par\Needspace{25mm}\vspace{6pt}\noindent\colorbox{black!12}{\makebox[7mm]{\bfseries\strut T}}\hspace{2mm}{\bfseries Probetest}\par\vspace{4pt}
\lsg{1.}{$u = \pi \cdot d$}{Die Borte liegt am Rand, gefragt ist der Umfang.}
\lsg{2.}{$A \approx 4{,}91$\,m²}{$r = 2{,}50 : 2 = 1{,}25$\,m; $A = \pi \cdot 1{,}25^2 \approx 4{,}91$\,m²}
\lsg{3.}{$r \approx 1{,}49$\,m}{$r^2 = 7 : \pi$; $r = \sqrt{7 : \pi} \approx 1{,}49$\,m}
\lsg{4.}{Stoff ja (2,54\,m²), Borte nein (5,65\,m)}{Stoff: Ja; $r = 0{,}9$\,m, $A = \pi \cdot 0{,}9^2 \approx 2{,}54$\,m² < 3\,m². Borte: Nein; $u = \pi \cdot 1{,}8 \approx 5{,}65$\,m > 5\,m, es fehlen $\approx 0{,}65$\,m}
\end{document}
